% R15 component; only input paths and enumerated location/grammar prose are edited.
% Source: revisions/2026-10-04-r15/manuscript/economic_scope.tex; commit 1cb3cc9efe135966ec228dfaf51c6841a6dfc97c.
% Generated by code/integrate.py; edit extraction rules there.
% Reviewed source: revisions/2026-10-04-r14/manuscript/economic_scope.tex, line 1, commit f5021cefa71492babfcbfa580e0c984f59a9de26
\section{Economic Applications of the Bellman Structure}\label{sec:economic-scope}
The motivating economic problem combines external choices with the costly adjustment of preferences. Recursive utility changes how a given policy is evaluated; preference adjustment adds an internal choice to its Hamiltonian; and strategic interaction changes whose deviations define improvement. These distinctions determine the economic interpretation of NBO. The capital economy studied above provides a setting in which policy-specific verification can be carried through to a high-dimensional diffusion. The preference, temporal-self, and game applications establish the corresponding economic objects and their own verification results.

\subsection{Costly preference adjustment and portfolio choice}
Let $u$ denote risk aversion and $X$ financial wealth. The household chooses consumption $c$, portfolio share $p$, and preference-adjustment effort $\theta$. The preference state is reflected on a declared interval $[\underline u,\bar u]\subset(1,\infty)$, and wealth follows the portfolio diffusion:
\begin{align}
 du&=\theta\,dt+\sigma_u\,dW^u+dL^{\underline u}-dL^{\bar u},\nonumber\\
 dX&=\{[r+p(\mu-r)]X-c\}\,dt+p\sigma_X X\,dW^X,
 \qquad d\langle W^u,W^X\rangle_t=\varrho\,dt.\label{eq:scope-preferences}
\end{align}
For a fixed consumption unit $\bar c>0$, flow utility is
\[
 U(c,u)-\tfrac12k\theta^2,
 \qquad U(c,u)=\frac{(c/\bar c)^{1-u}}{1-u}.
\]
The objective discounts both felicity and adjustment costs at rate $\rho$ and includes the specified wealth-exit and terminal payoff. The baseline uses $\bar c=1$. Because the household controls $u$, the cardinal normalization across preference states is a primitive: adding a function of $u$ changes the incentives to adjust preferences. The reflected state, positive consumption domain, and stopping data give the evaluation equation its economic and mathematical domain.

At an interior optimum with $V_X>0$ and $V_{XX}<0$, the three policy conditions are
\begin{align}
 c^*&=\bar c(\bar c V_X)^{-1/u}, &
 \theta^*&=V_u/k,\nonumber\\
 p^*&=-\frac{V_X(\mu-r)}{XV_{XX}\sigma_X^2}
       -\frac{\varrho\sigma_u V_{uX}}{XV_{XX}\sigma_X}.\label{eq:scope-preference-policies}
\end{align}
Box constraints require their corresponding feasible maximizers. The cross derivative in the portfolio rule links the external risk choice to the internal preference state. Its sign is an equilibrium object to be determined by evaluation. The adjustment rule likewise relates effort to the shadow value of changing preferences. A higher adjustment cost lowers the value of every fixed policy and hence the optimized value; whenever an envelope derivative exists,
\begin{equation}\label{eq:scope-preference-envelope}
 \partial_k V^k(t,u,X)
 =-\tfrac12\E^{a^k}_{t,u,X}
   \int_t^\tau e^{-\rho(s-t)}\theta_s^2\,ds\leq0.
\end{equation}
This comparative static holds with the feasible class and stopping payoff fixed.

The numerical application implements these choices jointly. The finite-action study evaluates complete policies by independent backward recursion. The continuous-action study trains the actor through its own payoff and verifies its hybrid policy with a complete action-box enclosure. Fresh training on nested grids separates the effect of resolution from the evaluation of a frozen coarse policy. A model-specific compensating-consumption calculation translates its cardinal utility bounds into an explicit transfer experiment. Application Sections~\ref{sec:ndu}, \ref{sec:ndu-neural}, and~\ref{sec:r9results} give the full model, all recorded policies, refinement evidence, and compensation account.

\subsection{Recursive utility and temporal selves}
Recursive utility enters policy evaluation through the policy's own continuation utility. With $a=1-1/\psi$, $b=1-\gamma$, and $bV>0$, the Epstein--Zin aggregator used here is
\begin{equation}\label{eq:scope-recursive}
 f(c,V)=\frac\rho a\{c^a(bV)^{1-a/b}-bV\}.
\end{equation}
The invariant utility interval and boundary payoff are part of the evaluation problem. Application Section~\ref{sec:recursive} derives the homothetic policy conditions and verifies the matching-bequest experiment, including the barriers that keep every policy utility in the aggregator's domain. This application tests recursive evaluation while maintaining the same separation between value approximation and feasible action improvement.

Time inconsistency requires a different improvement criterion. For the continuous discount kernel
\[
 h(s)=\beta e^{-\rho s}+(1-\beta)e^{-(\rho+\kappa)s},
 \qquad 0<\beta\leq1,\quad\kappa>0,
\]
two critics evaluate the same future policy at the two discount rates. Writing these values as $W_1,W_2$, the current self uses $V=\beta W_1+(1-\beta)W_2$ and chooses an instantaneous deviation by maximizing $u(a)+\LL^a V$. Application Section~\ref{sec:temporal} states the complete component equations and derives the spike-deviation condition. In the homothetic consumption economy, the equilibrium consumption share is the unique root in its admissible interval, and stronger present bias raises that share. Independently optimized finite-duration deviations check the limiting equilibrium condition. Thus the actor's objective follows from the temporal equilibrium concept, while the critics retain their policy-evaluation role.

\subsection{Capital investment and strategic deviations}
In the dynamic Cournot application, firm $i$ chooses investment in its own capital, with
\[
 dK^i=(I^i-\delta K^i)\,dt+\sigma_KK^i\,dW^i,
 \qquad q_i=AK^i.
\]
Its flow profit is $P_N(K)q_i-I^i-b_i(I^i)^2/2$, where the market-size normalization in $P_N$ is specified for each economy. Given rivals' policies, a critic evaluates the resulting own-firm payoff and the actor changes only its own action. Consequently, the relevant economic error is the gain from a complete unilateral deviation against the frozen rival profile.

Application Sections~\ref{sec:games} and~\ref{sec:dynamic-computation} specify the boundaries, investment sets, and market normalizations, and report independent full dynamic best responses to every principal finite neural profile. The resulting profit-unit bounds include all dates and states of the stated finite game. They also cover randomized and history-dependent unilateral deviations against those fixed Markov rivals, by finite dynamic programming. The record reports the cost of the exhaustive training oracle and any nonzero one-step equilibrium defects along with the final exploitability calculation.

These applications give the common method a precise economic interpretation. A recursive critic evaluates a utility process; preference and capital actors allocate resources and adjustment effort; a temporal actor satisfies a spike condition; and a strategic actor is assessed against unilateral best responses. Each application supplies the boundary conditions, approximation class, and error account required by its own economic question.
